\documentclass{beamer}
\usepackage{graphicx}
\usepackage{amsmath}
\usetheme{Warsaw}

\graphicspath{ {./images/} }

\title{Inversión de matrices}
\subtitle{Para la solución de sistemas de ecuaciones lineales}
\author{Alan Cerón Chávez}

\institute{FES Acatlán - UNAM}
\date{\today}


\begin{document}
    \begin{frame}
        \titlepage
    \end{frame}

    \begin{frame}
        \frametitle{Contenido}
        \tableofcontents[sections=1-2]
    \end{frame}

    \begin{frame}
        \frametitle{Contenido}
        \tableofcontents[sections=3-5]
    \end{frame}

    \section{Introducción}
    \subsection{¿Qué es una matriz?}

    \begin{frame}
        \frametitle{¿Qué es una matriz?}
        Una matriz $A$ es un arreglo rectangular de números de mn elementos, ordenados en m filas y n columnas.
        $$
        A=
            \begin{pmatrix}
                a_{11} & a_{12} & \cdots & a_{1n} \\
                a_{21} & a_{22} & \cdots & a_{2n} \\
                \vdots & \vdots & \ddots & \vdots \\
                a_{1m} & a_{2m} & \cdots & a_{nm} \\
            \end{pmatrix}
        $$
    \end{frame}


    \subsection{¿Cómo se relaciona con un sistema de ecuaciones lineales?}
    \begin{frame}
        \frametitle{Representación de un sistema de ecuaciones lineales como una matriz}
        sea el siguiente sistema de ecuaciones lineales:
        $$
        \begin{cases}
            2x + 3y = 8 \\
            3x + 4y = 9
        \end{cases}
        $$
        Los coeficientes de las variables se representan como una matriz:
        $$
        \begin{pmatrix}
            2 & 3 \\
            3 & 4
        \end{pmatrix}
        $$
        Y el sistema de ecuaciones lineales se puede representar de la manera $Ax=b$
        $$
        \begin{pmatrix}
            2 & 3 \\
            3 & 4
        \end{pmatrix}
        \begin{pmatrix}
            x \\
            y
        \end{pmatrix}
        =
        \begin{pmatrix}
            8 \\
            9
        \end{pmatrix}
        $$

    \end{frame}
    \subsection{Resolviendo el sistema de ecuaciones lineales}

    \begin{frame}
        Con la representación $Ax=b$ podemos resolver el sistema de ecuaciones lineales de la siguiente manera:
        \begin{align*}
            Ax &= b \\
            A^{-1}Ax &= A^{-1}b \\
            Ix &= A^{-1}b \\
            x &= A^{-1}b
        \end{align*}
    \end{frame}


    \section{Desarrollo}
    \subsection{¿Por qué?}

    \begin{frame}
        La matriz identidad $I$ es una matriz cuadrada que tiene unos en su diagonal principal y ceros en el resto de sus elementos.
        Tiene la propiedad de que al multiplicar cualquier matriz $A$ por $I$ el resultado es la misma matriz $A$.
        $$
        I =
        \begin{pmatrix}
            1 & 0 & \cdots & 0 \\
            0 & 1 & \cdots & 0 \\
            \vdots & \vdots & \ddots & \vdots \\
            0 & 0 & \cdots & 1 \\
        \end{pmatrix}
        $$

    \end{frame}


    \subsection{La inversa de una matriz}

    \begin{frame}
        \frametitle{Inversa de una matriz}
        El producto de una matriz $A$ por su inversa $A^{-1}$ es la matriz identidad $I$.
        Sea $A$ una matriz y su determinante es distinta de cero, entonces:
        $$
        A^{-1} = \frac{1}{det(A)}adj(A)
        $$
        $$
        A^{-1} = \frac{1}{det(A)}C^{T}
        $$
    \end{frame}

    \subsection{El detrminante de una matriz}

    \begin{frame}
        \frametitle{Caso 1: Matriz de 1x1}
        Sea $A$ una matriz de 1x1, entonces su determinante es el elemento de la matriz.
        $$
        A =
        \begin{pmatrix}
            a_{11}
        \end{pmatrix}
        \quad \det{A} = a_{11}
        $$
    \end{frame}

    \begin{frame}
        \frametitle{Caso 2: Matriz de 2x2}
        Sea $A$ una matriz de 2x2, entonces su determinante es:
        $$
        A =
        \begin{pmatrix}
            a_{11} & a_{12} \\
            a_{21} & a_{22}
        \end{pmatrix}
        \quad \det{A} = a_{11}a_{22} - a_{12}a_{21}
        $$
    \end{frame}

    \begin{frame}
        Sea $A$ una matriz de 3x3, entonces su determinante es:
        $$
        A =
        \begin{pmatrix}
            a_{11} & a_{12} & a_{13} \\
            a_{21} & a_{22} & a_{23} \\
            a_{31} & a_{32} & a_{33}
        \end{pmatrix}
        \quad \det{A} =
        \begin{align*}
            a_{11}(a_{22}a_{33} &- a_{23}a_{32}) \\
            &- \\
            a_{12}(a_{21}a_{33} &- a_{23}a_{31}) \\
            &+ \\
            a_{13}(a_{21}a_{32} &- a_{22}a_{31})
        \end{align*}
        $$
        O de la siguiente manera:
        $$
        \det{A} = a_{11}\begin{vmatrix}a_{22} & a_{23}\\ a_{32} & a_{33}\end{vmatrix} - a_{12}\begin{vmatrix}a_{21} & a_{23}\\ a_{31} & a_{33}\end{vmatrix} + a_{13}\begin{vmatrix}a_{21} & a_{22}\\ a_{31} & a_{32}\end{vmatrix}
        $$
    \end{frame}

        \begin{frame}
            \frametitle{Caso n: Matriz de nxn}
            Sea $A$ una matriz de nxn, entonces su determinante es:
            $$
            \det{A} = \sum_{i=1}^{n} a_{11}A_{11} + a_{12}A_{12} + \cdots + a_{1n}A_{1n}
            $$
        \end{frame}

    \subsection{Co-factores}

    \begin{frame}
        \begin{center}
            Qué es A{ij}?
        \end{center}\\
        A{ij} es llamado el cofactor de la matriz A en $ij$, y se define como:
        $$
        A_{ij} = (-1)^{i+j}\det{M_{ij}}
        $$

        Donde $M_{ij}$ es la matriz menor de $A$ en $ij$.\\
        Nota: $(-1)^{i+j}$ es el signo del cofactor. Y puede ser expresado como:
        $$
        (-1)^{i+j} =
        \begin{cases}
            1 & \text{si } i+j \text{ es par}\\
            -1 & \text{si } i+j \text{ es impar}
        \end{cases}
        $$
    \end{frame}

    \subsection{La matriz menor}

    \begin{frame}
        \frametitle{Qué es Mij?}
        Sea $A$ una matriz de nxn, entonces su matriz menor $M_{ij}$ de dimensiones (n-1)x(n-1) es la matriz resultante de eliminar la fila $i$ y la columna $j$ de la matriz $A$.\\
        Ejemplo:
        $$
        A =
        \begin{pmatrix}
            1 & 2 & 3 \\
            4 & 5 & 6 \\
            7 & 8 & 9
        \end{pmatrix}
        \quad M_{12} =
        \begin{pmatrix}
            4 & 6 \\
            7 & 9
        \end{pmatrix}
        $$
    \end{frame}

    \subsection{Matriz de cofactores}

    \begin{frame}
        \frametitle{Matriz de cofactores}
        Con todo lo anterior podemos definir la matriz de cofactores y la adjunta de una matriz.\\
        Comenzando con la matriz de cofactores:\\
        La Matriz de cofactores $C$ de una matriz $A$ es la matriz que contiene los cofactores de $A$.
        $$
        C = \begin{pmatrix}
            A_{11} & A_{12} & \cdots & A_{1n} \\
            A_{21} & A_{22} & \cdots & A_{2n} \\
            \vdots & \vdots & \ddots & \vdots \\
            A_{n1} & A_{n2} & \cdots & A_{nn}
        \end{pmatrix}
        $$
    \end{frame}

    \begin{frame}
        \frametitle{Matriz adjunta}
        La matriz adjunta es la matriz de cofactores transpuesta.
        $$
        adj(A) = C^T
        $$
    \end{frame}

    \subsection{Inversión de matrices}

    \begin{frame}
        \frametitle{Inversión de matrices}
        Finalmente, podemos definir la inversa de una matriz.\\
        Regresamos a la definición del inicio:
        $$
        A^{-1} = \frac{1}{\det{A}}adj(A)
        $$
        Y Si tenemos una matriz de $4x4$? Cuantas operaciones tenemos que hacer?\\
        \begin{center}
            ¿Para $n x n$?
        \end{center}
    \end{frame}
    \section{Ejercicios}
    \subsection{Ejercicio 1}

    \begin{frame}
        \frametitle{Matriz de 4x4}
        Sea la matriz $A$ de 4x4:
        $$
        A =
        %nmms copilot haz algo más complejo
        % \begin{pmatrix}
        %     1 & 2 & 3 & 4 \\
        %     5 & 6 & 7 & 8 \\
        %     9 & 10 & 11 & 12 \\
        %     13 & 14 & 15 & 16
        % \end{pmatrix}
        \begin{pmatrix}
            25 & 36 & 57 & -12 \\
            54 & 23 & 12 & 45 \\
            12 & 45 & 78 & 96 \\
            45 & 78 & 96 & 12
        \end{pmatrix}
        $$
        Calcular la inversa de $A$.\\
        $$
        A^{-1} = \frac{1}{\det{A}}adj(A)
        $$
    \end{frame}

    \subsection{Ejercicio 2}

    \begin{frame}
        \frametitle{Matriz de 8x8}
        Sea la matriz $A$ de 8x8:
        $$
        A =
        %Matriz de 8x8 con determinante diferente de 0
        \begin{pmatrix}
              3 & 2 & 1 & 4 & 0 & 5 & 6 & 7 \\
              0 & 8 & 9 & 0 & 1 & 0 & 2 & 0 \\
              7 & 0 & 3 & 0 & 6 & 2 & 0 & 8 \\
              4 & 0 & 5 & 9 & 0 & 1 & 0 & 0 \\
              0 & 6 & 0 & 0 & 7 & 0 & 9 & 0 \\
              1 & 0 & 0 & 0 & 8 & 4 & 0 & 3 \\
              0 & 0 & 4 & 5 & 2 & 0 & 7 & 1 \\
              2 & 0 & 0 & 3 & 0 & 9 & 8 & 6 \\
        \end{pmatrix}
        $$
        Calcular la inversa de $A$.\\
        $$
        A^{-1} = \frac{1}{\det{A}}adj(A)
        $$
    \end{frame}
    \subsection{Ejercicio 3}
        \begin{frame}
            \frametitle{Resolver el sistema de ecuaciones}
            Sea el sistema de ecuaciones:
            $$
            \begin{cases}
                12x + 3y - 4z = 17 \\
                -3x + 24y + 45z = -6 \\
                x + 9y - 6z = -13
            \end{cases}
            $$
            Resolver el sistema de ecuaciones utilizando la matriz inversa.
            $$
                A^{-1} = \frac{1}{\det{A}}adj(A)\qquad x = A^{-1}b
            $$
        \end{frame}

    \section{implementación}
    \subsection{QR al código}
    \begin{frame}
        \frametitle{Implementación en Rust}
        \begin{center}
            \includegraphics[width=170pt]{images/QR_rust_implementation.png}
        \end{center}
    \end{frame}
    \subsection{Live en Repl.it}
    \section{Conclusiones}
    \subsection{Complejidad computacional}

    \begin{frame}
        \frametitle{Complejidad computacional}
        \begin{itemize}
            \item El cálculo de la inversa de una matriz es un proceso que requiere de mucho tiempo de cómputo.\\
            \item La complejidad computacional de este proceso es de $O(n^3)$\\
            \item Por lo que es necesario implementar algoritmos que nos permitan obtener la inversa de una matriz de manera eficiente.\\
        \end{itemize}
    \end{frame}

    \subsection{Conclusiones}
    \begin{frame}
        \frametitle{Conclusiones}
        \begin{itemize}
            \item El cálculo de la inversa de una matriz nos permite resolver sistemas de ecuaciones lineales.\\
            \item Su obtención es un proceso tedioso y que requiere de mucho tiempo de cómputo.\\
            \item Por lo que es necesario implementar algoritmos que nos permitan obtener la inversa de una matriz de manera eficiente.\\
            \item Aunque existen metodos para resolver sistemas de ecuaciones lineales que no requieren de la inversa de una matriz.\\
            \item En las próximas presentaciones se verán algunos de estos métodos.
        \end{itemize}
    \end{frame}

\end{document}
